% !TeX program = lualatex
% !TeX encoding = utf8
% !TeX spellcheck = uk_UA
% !TeX root = ../ElMagConspect.tex
% УВАГА: для \enquote{...} потрібен пакет csquotes (\usepackage[ukrainian=quotes]{csquotes}
% в преамбулі ElMagConspect.tex) --- переконайтесь, що він підключений.

%=========================================================
\chapter{Енергія електричного поля}\label{\currfilebase}
\graphicspath{{\currfilebase/Pictures}}
%=========================================================
\epigraph{Енергія електричного поля --- це не властивість зарядів, а властивість простору навколо них: поле не лише переносить взаємодію між
	зарядами, а й саме є носієм і сховищем енергії.
}{Переосмислення ідей сучасної фізики}

У попередніх розділах ми розібрали, як поводиться електричне поле у вакуумі, у провідниках і в діелектриках. Ми знаємо, як обчислити поле за
заданим розподілом зарядів (теорема Гаусса, принцип суперпозиції), і як поводиться поле на межі розділу середовищ. Тепер настав час поставити
наступне природне запитання: \emphx{скільки енергії зберігає електричне поле} і \emphx{як цю енергію використати для обчислення сил}, що діють на
заряджені тіла?

Цей розділ складається з трьох частин. Спершу ми введемо \emphx{електричну ємність} --- величину, яка описує, скільки заряду (а отже, і поля)
здатна накопичити система провідників при заданій різниці потенціалів. Далі ми покажемо, що енергія заряджених тіл --- це насправді енергія
\emphx{поля}, розподілена в просторі з густиною $w = \Efield\cdot\Dfield/8\pi$. Нарешті, ми навчимось обчислювати \emphx{пондеромоторні сили} ---
сили, що діють на заряджені провідники й діелектрики в електричному полі, --- за допомогою методу віртуальних переміщень.

%% --------------------------------------------------------
\section{Електрична ємність}
%% --------------------------------------------------------

\subsection*{Ємність відокремленого провідника}

Нехай усамітнений (далекий від інших тіл) провідник несе заряд $q$; тоді, як ми знаємо, весь цей заряд розподіляється по його поверхні, а сам
провідник --- еквіпотенціальне тіло з деяким потенціалом $\phi$. Якщо заряд збільшити в $n$ разів, то, за принципом суперпозиції, поле в кожній точці
простору також зросте в $n$ разів, а отже, в $n$ разів зросте й робота з переміщення пробного заряду з поверхні провідника на нескінченність, тобто й
сам потенціал $\phi$. Це означає, що відношення
\begin{equation}\label{eq:CapacitorDef}
	\tcbhighmath{
		C = \frac{q}{\phi}
	}
\end{equation}
не залежить від заряду провідника і визначається лише його формою, розмірами та властивостями середовища, що його оточує. Цю величину називають
\emphx{електричною ємністю провідника}.

\begin{Attention}
	Ємність провідника --- суто геометрична характеристика: вона залежить від форми й розмірів тіла, а також від діелектричної проникності
	середовища, але \emphx{не залежить} ні від заряду, ні від матеріалу самого провідника. Означення~\eqref{eq:CapacitorDef} коректне лише для
	\emphx{усамітненого} провідника, тобто коли поблизу немає інших заряджених тіл, здатних істотно змінити його потенціал.
\end{Attention}

%% --------------------------------------------------------
\subsection*{Коефіцієнти ємності}
%% --------------------------------------------------------

Для усамітненого провідника заряд і потенціал пропорційні: \(q = C\phi\). Якщо провідників кілька, потенціал кожного залежить від зарядів усіх сусідів (принцип суперпозиції), тому
\begin{equation}\label{eq:Cij}
	q_i = \sum_j C_{ij}\,\phi_j.
\end{equation}
Коефіцієнти \(C_{ij}\) залежать лише від форми, розмірів і взаємного розташування провідників (і від проникності середовища), але не від самих зарядів чи потенціалів.

Діагональні \(C_{ii}\) --- \emphx{власні ємності}: заряд провідника \(i\) при одиничному власному потенціалі й заземлених сусідах. Недіагональні \(C_{ij}\) (\(i\neq j\)) --- \emphx{коефіцієнти взаємного впливу}: заряд, наведений на заземленому провіднику \(i\), коли сусід \(j\) піднято до одиничного потенціалу.

\begin{Attention}
	\emphx{Знаки.} \(C_{ii}>0\), а \(C_{ij}\le 0\) при \(i\neq j\): піднімаючи потенціал одного провідника, ми завжди наводимо на сусідніх заряд \emphx{протилежного} знака.

	\emphx{Симетрія.} \(C_{ij}=C_{ji}\) --- наслідок теореми взаємності електростатики. Порівняємо два стани: (1) піднято до потенціалу \(1\) провідник \(j\), решта заземлені, тоді \(q_i^{(1)}=C_{ij}\); (2) навпаки, піднято провідник \(i\), тоді \(q_j^{(2)}=C_{ji}\). Теорема взаємності \(\sum_k q_k^{(1)}\phi_k^{(2)} = \sum_k q_k^{(2)}\phi_k^{(1)}\) при єдиному ненульовому доданку в кожній сумі дає \(C_{ij}=C_{ji}\).
\end{Attention}

\emphx{Оцінка зверху для \(C_{12}\).} Розглянемо пару провідників з однаковою власною ємністю \(C_{11}=C_{22}\equiv C_0\) (наприклад, дві однакові кулі). Піднімемо потенціал провідника~1 до одиниці, а провідник~2 заземлимо (\(\phi_2=0\)). Тоді, згідно з~\eqref{eq:Cij}, заряд першого провідника дорівнює \(q_1=C_{11}=C_0\), а на другому наводиться заряд \(q_2=C_{12}\) (від'ємний, за щойно сформульованим правилом знаків).

	\begin{SCfigure}[0.75][h!]\centering
		\localinput{field-lines-bound.tikz}
		\caption{Частина силових ліній провідника~1 замикається на заземленому провіднику~2, решта йде на нескінченність\label{pic:FieldLinesBound}}
	\end{SCfigure}

Фізично \(q_2\) --- це заряд, \enquote{перехоплений} тими силовими лініями поля провідника~1, що замикаються саме на заземленому сусідові (рис.~\ref{pic:FieldLinesBound}); решта ліній іде на нескінченність (або на інші тіла). Оскільки сусід не може перехопити більше силових ліній, ніж їх випустив сам провідник~1, наведений заряд не може за модулем перевищити початковий:
\begin{equation}\label{eq:OffDiagonalBound}
	\tcbhighmath{
		|C_{12}| \le C_0.
	}
\end{equation}
Рівність \(C_{12}\to -C_0\) означає, що \emph{всі} силові лінії від провідника~1 замикаються на провіднику~2 --- жодна не йде на нескінченність. Саме до цієї межі прямує система, що перетворюється на конденсатор.

%% --------------------------------------------------------
\subsection*{Два провідники: взаємна ємність}
%% --------------------------------------------------------

Найпоширеніший спосіб зарядити провідник --- перенести частину заряду з другого: тоді на одному з'являється заряд \(+q\), на іншому \(-q\). Для такого --- дуже поширеного на практиці --- стану зручно ввести окрему характеристику пари замість громіздкої матриці \(C_{ij}\):
\begin{equation}\label{eq:MutualCapacitanceDef}
	\tcbhighmath{
		C = \frac{q}{\phi_+ - \phi_-}
	}
\end{equation}
де \(\phi_\pm\) --- потенціали провідників із зарядами \(\pm q\). Величину~\eqref{eq:MutualCapacitanceDef} називають \emphx{взаємною ємністю} пари.

Щоб виразити \(C\) через коефіцієнти \(C_{ij}\), підставимо \(q_1=+q\), \(q_2=-q\) у систему~\eqref{eq:Cij} (з урахуванням симетрії \(C_{21}=C_{12}\)):
\begin{equation*}
	q = C_{11}\phi_1 + C_{12}\phi_2, \qquad
	-q = C_{12}\phi_1 + C_{22}\phi_2.
\end{equation*}
Розв'язавши цю лінійну систему відносно \(\phi_1\), \(\phi_2\) (звичайним оберненням \(2\times2\)-матриці), знаходимо, що різниця потенціалів пропорційна заряду:
\begin{equation*}
	\phi_1 - \phi_2 = \frac{q\,(C_{11}+C_{22}+2C_{12})}{C_{11}C_{22}-C_{12}^2},
\end{equation*}
звідки за означенням~\eqref{eq:MutualCapacitanceDef}
\begin{equation}\label{eq:MutualCapacitanceGeneral}
	\tcbhighmath{
		C = \frac{C_{11}C_{22}-C_{12}^2}{C_{11}+C_{22}+2C_{12}}.
	}
\end{equation}

Для \emphx{симетричної пари} \(C_{11}=C_{22}\equiv C_0\) чисельник розкладається як різниця квадратів, \(C_0^2-C_{12}^2=(C_0-C_{12})(C_0+C_{12})\), а знаменник дорівнює \(2(C_0+C_{12})\); після скорочення спільного множника \((C_0+C_{12})\):
\begin{equation}\label{eq:SymmetricPairCap}
	\tcbhighmath{
		C = \frac{C_0 - C_{12}}{2}.
	}
\end{equation}

Ця формула добре узгоджується з оцінкою~\eqref{eq:OffDiagonalBound} і показує, як взаємна ємність зростає в міру екранування: якщо взаємний вплив зникає (\(C_{12}\to 0\), провідники далеко один від одного), то \(C\to C_0/2\) --- це просто послідовне з'єднання двох однакових ємностей \(C_0\); якщо ж екранування майже повне (\(C_{12}\to -C_0\), гранично допустиме за~\eqref{eq:OffDiagonalBound}), то \(C\to C_0\) --- поле майже все замкнене між обкладками. Саме другий режим реалізується в конденсаторі.

\emphx{Приклад: дві далекі кулі.} Дві кулі на відстані \(b\), значно більшій за їхні розміри (рис.~\ref{pic:TwoSpheres}). Нехай \(C_1\), \(C_2\) --- ємності куль поодинці (\(C_{11}\to C_1\), \(C_{22}\to C_2\) при \(b\to\infty\)).

	\begin{SCfigure}[1.5][h!]\centering
		\localinput{two-spheres.tikz}
		\caption{Дві далекі одна від одної заряджені кулі\label{pic:TwoSpheres}}
	\end{SCfigure}

	Зарядимо кулі \(\pm q\) і врахуємо вплив сусідньої кулі як точкового заряду на відстані \(b\):
	\begin{equation*}
		\phi_+ = \frac{q}{C_1} - \frac{q}{b}, \qquad
		\phi_- = -\frac{q}{C_2} + \frac{q}{b}.
	\end{equation*}
	Звідси \(\phi_+-\phi_- = q\left(\dfrac1{C_1}+\dfrac1{C_2}-\dfrac2b\right)\), тому
	\begin{equation*}
		C = \frac{C_1 C_2}{C_1 + C_2 - \dfrac{2C_1C_2}{b}}
		\approx \frac{C_1C_2}{C_1+C_2}\left(1 + \frac{C_1C_2}{C_1+C_2}\frac{2}{b}\right).
	\end{equation*}
	\begin{equation}\label{eq:TwoSpheresCap}
		\tcbhighmath{
			C \approx \frac{C_1C_2}{C_1+C_2}\left(1 + \frac{C_1C_2}{C_1+C_2}\frac{2}{b}\right).
		}
	\end{equation}
	У головному наближенні (\(b\to\infty\)) взаємна ємність дорівнює ємності послідовного з'єднання; поправка описує слабкий взаємний вплив куль.

	Той самий результат випливає й з коефіцієнтів \(C_{ij}\): якщо провідник~2 заземлити, для утримання \(\phi_2=0\) на ньому має навестись заряд \(q_2\approx -C_2q_1/b\) (компенсація поля точкового заряду \(q_1\) на відстані \(b\)), звідки \(C_{12}\approx -C_1C_2/b\) --- від'ємний, як і належить за загальним правилом знаків. Підстановка в~\eqref{eq:MutualCapacitanceGeneral} відтворює~\eqref{eq:TwoSpheresCap}: обидва підходи узгоджені, і саме знак \(C_{12}\) тут принциповий --- зворотний знак дав би суперечливу відповідь.

%% --------------------------------------------------------
\subsection*{Конденсатор}
%% --------------------------------------------------------

\emphx{Конденсатор} --- система з двох провідників (\emphx{обкладок}), розташованих так близько (зазвичай через тонкий шар діелектрика), що майже все поле зосереджене у проміжку між ними. Завдяки такому \emphx{екрануванню} взаємна ємність конденсатора, на відміну від ємності одинокого провідника чи пари віддалених провідників, практично не залежить від інших тіл поблизу --- саме тому вона є зручною й відтворюваною характеристикою. Формула~\eqref{eq:MutualCapacitanceDef} лишається тією самою; конденсатор --- окремий випадок пари провідників з екранованим полем.

Наскільки повним може бути це екранування, видно з~\eqref{eq:SymmetricPairCap}: коли обкладки зближуються, коефіцієнт взаємного впливу прямує до своєї межі \(C_{12}\to -C_0\) (оцінка~\eqref{eq:OffDiagonalBound}), а отже,
\begin{equation*}
	C \to C_0,
\end{equation*}
тобто взаємна ємність прямує до власної. У цій границі поле майже все замкнене між обкладками, і \(C\) перестає залежати від того, що відбувається далеко.

%% --------------------------------------------------------
\subsection*{З'єднання конденсаторів}
%% --------------------------------------------------------

На практиці конденсатори часто з'єднують у батареї --- послідовно або паралельно, --- щоб отримати потрібну сумарну ємність.

\begin{figure}[h!]
	\centering
	% --- (а) послідовне з'єднання ---
	\begin{subfigure}[t]{0.48\linewidth}
		\centering
		\begin{circuitikz}[>=latex, scale=1.5]
			% Послідовне з'єднання конденсаторів
			\draw (0,0)
			node[xshift=5pt, above, font=\tiny, text=red] {$+q$}
			to[C=$C_1$]
			node[above, font=\tiny, text=blue] {$-q$}
			++(1,0)
			node[xshift=5pt, above, font=\tiny, text=red] {$+q$}
			to[C=$C_2$]
			node[above, font=\tiny, text=blue] {$-q$}
			++(1,0)
			node[xshift=5pt, above, font=\tiny, text=red] {$+q$}
			to[C=$C_3$]
			node[above, font=\tiny, text=blue] {$-q$}
			++(1,0)
			node[right] {\ldots} ++(0.5,0)
			node[xshift=5pt, above, font=\tiny, text=red] {$+q$}
			to[C=$C_n$]
			node[above, font=\tiny, text=blue] {$-q$}
			++(1,0) coordinate (C) node[circle, fill, inner sep=1pt] {} ;
			\draw (0,0) node[circle, fill, inner sep=1pt] {} -- ++(0, -1);
			\draw (C) -- ++(0, -1);
			\draw[<->] ([yshift=-0.75cm]0,0)  -- node[below,
				font=\scriptsize] {$U = U_1+U_2
					+U_3 + \ldots +
					U_n$}
			([yshift=-0.75cm]C);
		\end{circuitikz}
		\caption{Послідовне з'єднання}
		\label{pic:CapacitorsSeries}
	\end{subfigure}\hfill
	% --- (б) паралельне з'єднання ---
	\begin{subfigure}[t]{0.48\linewidth}
		\centering
		\begin{circuitikz}[>=latex, scale=1, transform shape, rotate=90]
%         \ctikzset{label/align=rotate}
			% Верхня лінія
			\foreach \y in {1,...,5} {
					\ifnum\y=4
						\node at (1,{1.5*\y}) {$\ldots$};
					\else
						\ifnum\y=5
							\edef\c{n}
						\else
							\edef\c{\y}
						\fi
						\draw (0, {1.5*\y}) to[C, l={\rotatebox[origin=c]{-90}{$C_\c$}}] ++(2,0);
					\fi

				}
			\draw[red, thick] (0,1.5) -- ++(0,{4*1.5});
			\draw[blue, thick]  (2,1.5) -- ++(0,{4*1.5});

			\draw[thick, red] (-1, {3*1.5}) node[circle, fill, inner sep=1pt] (A) {} -- ++(1,
			0) ;
			\draw[thick, blue] (2, {3*1.5})   -- ++(1, 0) node[circle, fill,
				inner sep=1pt] (B) {} ;

			\draw[] (A) -- ++(0, -4) coordinate (D) ;
			\draw[] (B) -- ++(0, -4) coordinate (E);
			\draw[<->] ([yshift=0.25cm]D) -- node[right, rotate=-90] {$U$} ([yshift=0.25cm]E);

			\fill[opacity=0.2, dashed, red!40,  very thin] (-0.55,1) rectangle ++(1.5, 7);
			\fill[opacity=0.2, dashed, blue!40, very thin] (1.05,1) rectangle ++(1.5, 7);
		\end{circuitikz}
		\caption{Паралельне з'єднання}
		\label{pic:CapacitorsParallel}
	\end{subfigure}
	\caption{З'єднання конденсаторів: (а) послідовне --- заряд на всіх обкладках однаковий, напруги додаються;
		(б) паралельне --- напруга на всіх конденсаторах однакова, заряди додаються}
	\label{pic:CapacitorsConnection}
\end{figure}

\emphx{Послідовне з'єднання.} Нехай $n$ конденсаторів з'єднано послідовно (рис.~\ref{pic:CapacitorsSeries}). Оскільки обкладки сусідніх
конденсаторів, з'єднані між собою провідником, електрично ізольовані від решти кола, повний заряд кожної внутрішньої пари обкладок мусить
залишатися нульовим: якщо на крайню ліву обкладку батареї нанести заряд $+q$, то шляхом електростатичної індукції той самий за модулем заряд $q$
з'явиться на \emphx{кожному} конденсаторі батареї. Загальна різниця потенціалів (напруга) на батареї дорівнює сумі напруг на окремих конденсаторах:
\begin{equation*}
	U = U_1 + U_2 + \ldots + U_n = \frac{q}{C_1} + \frac{q}{C_2} + \ldots + \frac{q}{C_n}.
\end{equation*}
Оскільки ємність батареї за означенням $C = q/U$, маємо
\begin{equation}\label{eq:SeriesCap}
	\tcbhighmath{
		\frac1C = \sum\limits_{i=1}^n \frac1{C_i}.
	}
\end{equation}

\emphx{Паралельне з'єднання.} Нехай тепер $n$ конденсаторів з'єднано паралельно (рис.~\ref{pic:CapacitorsParallel}): усі ліві обкладки з'єднані між
собою і мають спільний потенціал, так само як і всі праві. Тому напруга на кожному конденсаторі однакова і дорівнює загальній напрузі батареї $U$, а
заряди на окремих конденсаторах додаються:
\begin{equation*}
	q = q_1 + q_2 + \ldots + q_n = C_1 U + C_2 U + \ldots + C_n U.
\end{equation*}
Звідси
\begin{equation}\label{eq:ParallelCap}
	\tcbhighmath{
		C = \sum\limits_{i=1}^n C_i.
	}
\end{equation}

\begin{Attention}
	Послідовне з'єднання завжди \emphx{зменшує} сумарну ємність порівняно з найменшою з ємностей $C_i$ (аналогія з опорами резисторів у
	паралельному з'єднанні), тоді як паралельне з'єднання завжди \emphx{збільшує} сумарну ємність. Це протилежно до того, як поводяться опори при
	послідовному й паралельному з'єднанні резисторів, і пов'язано з тим, що ємність --- це \enquote{здатність накопичувати заряд}, а не \enquote{перешкода струму}.
\end{Attention}

\subsection*{Обчислення ємності найпростіших конденсаторів}

Усі задачі на обчислення ємності розв'язуються за \emphx{єдиним алгоритмом}:

\begin{enumerate}
	\item Задати заряд $q$ на обкладках і знайти розподіл поля $\Efield$ між ними (за теоремою Гаусса, з урахуванням діелектрика, якщо він є).
	\item Обчислити різницю потенціалів $\Delta\phi = \phi_+ - \phi_- = \int \Efield\cdot d\vect\ell$ уздовж силової лінії.
	\item За означенням~\eqref{eq:MutualCapacitanceDef} знайти $C = q/\Delta\phi$.
\end{enumerate}

Продемонструємо цей алгоритм на кількох типових прикладах. Скрізь вважатимемо, що простір між обкладками заповнений однорідним діелектриком із
проникністю $\epsilon$ (результати про поле в діелектрику --- з розділу~\ref{FileldInDielectrics}).

\emphx{Ємність усамітненої кулі.} Провідна куля радіусом $R$, що несе заряд $Q$, створює поза собою поле точкового заряду, а її власний потенціал
дорівнює потенціалу на поверхні, $\phi = Q/(\epsilon R)$. Звідси відразу
\begin{equation}\label{eq:SphereCap}
	\tcbhighmath{
		C = \epsilon R.
	}
\end{equation}
Цей простий результат корисно пам'ятати як орієнтир: ємність провідного тіла має порядок величини його лінійних розмірів (помножених на $\epsilon$).

\emphx{Плоский конденсатор.} Нехай обкладки --- дві паралельні пластини площею $S$ кожна, розташовані на відстані $d \ll \sqrt{S}$ одна від одної
(щоб знехтувати крайовими ефектами), а простір між ними заповнений діелектриком з проникністю $\epsilon$. Поле між обкладками однорідне,
$E = 4\pi\sigma/\epsilon = 4\pi q /(\epsilon S)$, тому різниця потенціалів $\Delta\phi = Ed = 4\pi qd/(\epsilon S)$, і
\begin{equation}\label{eq:FlatCap}
	\tcbhighmath{
		C = \frac{\epsilon S}{4\pi d}.
	}
\end{equation}

\emphx{Сферичний конденсатор.} Обкладки --- дві концентричні сфери радіусів $R_1 < R_2$, простір між якими заповнений діелектриком з проникністю
$\epsilon$. За теоремою Гаусса поле між обкладками таке саме, як і поле точкового заряду $q$ у діелектрику, $E = q/(\epsilon r^2)$. Різниця
потенціалів
\begin{equation*}
	\Delta\phi = \int\limits_{R_1}^{R_2} E\, dr = \frac{q}{\epsilon}\left(\frac1{R_1} - \frac1{R_2}\right) = \frac{q}{\epsilon}\frac{R_2 - R_1}{R_1
		R_2},
\end{equation*}
звідки
\begin{equation}\label{eq:SphericalCap}
	\tcbhighmath{
		C = \frac{\epsilon R_1 R_2}{R_2 - R_1}.
	}
\end{equation}
Легко переконатись, що при $R_2 \to \infty$ (друга обкладка йде на нескінченність) формула~\eqref{eq:SphericalCap} переходить у
формулу~\eqref{eq:SphereCap} для усамітненої кулі, як і повинно бути.

\emphx{Циліндричний конденсатор.} Обкладки --- дві коаксіальні циліндричні поверхні радіусів $R_1 < R_2$ і довжини $\ell \gg R_2$ (щоб знехтувати
крайовими ефектами), простір між ними заповнений діелектриком з проникністю $\epsilon$. За теоремою Гаусса поле між обкладками
$E = \dfrac{2q}{\epsilon \ell r}$ (поле нескінченної зарядженої нитки в діелектрику), тому
\begin{equation*}
	\Delta\phi = \int\limits_{R_1}^{R_2} E\, dr = \frac{2q}{\epsilon\ell}\ln\frac{R_2}{R_1},
\end{equation*}
звідки
\begin{equation}\label{eq:CylindricalCap}
	\tcbhighmath{
		C = \frac{\epsilon \ell}{2\ln\left(R_2/R_1\right)}.
	}
\end{equation}

\begin{Attention}
	Зверніть увагу, як \emphx{однаково} влаштовані всі чотири формули: чисельник --- це добуток $\epsilon$ на характерний \enquote{геометричний
	розмір} (площа / відстань, або добуток радіусів / різниця радіусів, або довжина / логарифм відношення), а знаменник --- безрозмірний
	геометричний множник. Запам'ятовувати окремо кожну формулу не потрібно --- достатньо вміти вивести її за алгоритмом вище.
\end{Attention}

\begin{Example}{Числовий приклад}
	Розглянемо плоский конденсатор із площею обкладок $S = 100\ \text{см}^2 = 10^{-2}\ \text{м}^2$ і відстанню між ними
	$d = 1\ \text{мм} = 10^{-3}\ \text{м}$, заповнений повітрям ($\epsilon \approx 1$). У гаусовій системі одиниць ємність такого конденсатора
	$C = \epsilon S/(4\pi d) = 10^{-2}/(4\pi\cdot 10^{-3})\ \text{см} \approx 0{,}8\ \text{см}$. У практичних одиницях (СІ) $1\ \text{см}$ ємності
	відповідає приблизно $1{,}1\ \text{пФ}$, тож $C \approx 0{,}9\ \text{пФ}$. Якщо між обкладками вставити слюду з $\epsilon \approx 6$, ємність
	зросте до $\approx 5\ \text{пФ}$ --- саме так \enquote{працює} діелектрик у конденсаторі: не змінюючи геометрії, він у $\epsilon$ разів збільшує
	ємність.
\end{Example}

\begin{Attention}
	\emphx{Суперконденсатори} (ультраконденсатори) мають ємність у тисячі фарад --- у мільярди разів більшу за ємність Землі
	($\approx 710\ \text{мкФ}$ як усамітненого провідника). Такої ємності досягають завдяки двом факторам:
	\begin{itemize}
		\item \emph{надрозвинена поверхня} пористих вуглецевих електродів ($1000$--$2600\ \text{м}^2/\text{г}$);
		\item \emph{подвійний електричний шар} --- іони електроліту розділені відстанню менше $1\ \text{нм}$.
	\end{itemize}
	Але навіть така ємність не робить їх універсальним накопичувачем. Порівняймо питому енергію:
	\begin{itemize}
		\item \emph{звичайний конденсатор}: $\sim 0{,}01$--$0{,}1\ \text{Вт·год/кг}$;
		\item \emph{суперконденсатор}: $\sim 5$--$10\ \text{Вт·год/кг}$ (до $30$ у лабораторних зразках);
		\item \emph{літій-іонний акумулятор}: $\sim 100$--$250\ \text{Вт·год/кг}$ --- енергія запасається в хімічних зв'язках, а не в полі;
		\item \emph{бензин}: $\approx 12\,000\ \text{Вт·год/кг}$.
	\end{itemize}
	Тому суперконденсатори не \emphx{витісняють} акумулятори, а доповнюють їх там, де важлива питома \emphx{потужність}, а не сумарна запасена енергія: регенеративне гальмування, стабілізація напруги, пікові навантаження.
\end{Attention}

%% --------------------------------------------------------
\section{Енергія електричного поля}
%% --------------------------------------------------------

\subsection*{Взаємна енергія системи точкових зарядів}

Щоб зблизити два точкових заряди $q_1$ і $q_2$ з нескінченності до відстані $r_{12}$ один від одного, потрібно виконати проти сил поля роботу
$A = q_1q_2/r_{12}$. Цю роботу і ототожнюють з \emphx{взаємною енергією} пари зарядів: $U_{12} = q_1q_2/r_{12}$. Для системи з $n$ точкових зарядів
(рис.~\ref{pic:ChargesTriangle}) сумарна взаємна енергія набирається з енергій усіх можливих пар:
\begin{equation}\label{eq:PairEnergy}
	\tcbhighmath{
		U = \frac12 \sum\limits_{i=1}^n \sum\limits_{\substack{j=1\\ j\neq i}}^n \frac{q_iq_j}{r_{ij}},
	}
\end{equation}
де множник $1/2$ компенсує подвійний перелік кожної пари (доданок з індексами $i,j$ і доданок з індексами $j,i$ описують одну й ту саму пару).

\begin{SCfigure}[1.5][h!]\centering
	\localinput{charges-triangle.tikz}
	\caption{Система трьох точкових зарядів: взаємна енергія набирається із внесків усіх пар\label{pic:ChargesTriangle}}
\end{SCfigure}

Формулу~\eqref{eq:PairEnergy} зручно переписати через потенціал. Потенціал поля в точці розташування $i$-го заряду, створюваний \emphx{усіма іншими}
зарядами (крім самого $i$-го, поле якого в точці свого ж розташування не визначене),
\begin{equation*}
	\phi_i = \sum\limits_{\substack{j=1 \\ j \neq i}}^n \frac{q_j}{r_{ij}},
\end{equation*}
тому
\begin{equation}\label{eq:SystemEnergyViaPotential}
	\tcbhighmath{
		W = \frac12 \sum\limits_i q_i \phi_i.
	}
\end{equation}

\begin{Attention}
	Формула~\eqref{eq:PairEnergy} описує \emphx{лише взаємну} енергію --- роботу, потрібну для зближення \emphx{вже готових} точкових зарядів з
	нескінченності. Вона \emphx{не враховує} власну енергію кожного точкового заряду окремо (докладніше про це --- наприкінці цього розділу). Для
	системи зарядів, що мають скінченні розміри, формула~\eqref{eq:PairEnergy} --- наближена, і тим точніша, чим менші розміри заряджених тіл
	порівняно з відстанями між ними.
\end{Attention}

\subsection*{Електростатична енергія тіла з неперервним розподілом заряду}

Для тіл з неперервним розподілом заряду --- об'ємною густиною $\rho(\vect r)$ та (якщо є) поверхневою густиною $\sigma(\vect r)$ ---
формулу~\eqref{eq:SystemEnergyViaPotential} узагальнюють природним чином, замінюючи суму на інтеграл, а $q_i$ --- на елемент заряду $\rho\, dV$ або
$\sigma\, dS$:
\begin{equation}\label{eq:BodyEnergy}
	\tcbhighmath{
		W = \frac12 \iiint\limits_V \rho(\vect r)\phi(\vect r)\, dV + \frac12 \iint\limits_S \sigma(\vect r)\phi(\vect r)\, dS.
	}
\end{equation}

\emphx{Приклад 1.} Нехай куля радіуса $R$ несе повний заряд $Q$, рівномірно розподілений по \emphx{поверхні}. Потенціал усюди на поверхні однаковий і
дорівнює $\phi = Q/R$, тому
\begin{equation}\label{eq:SphereEnergySurface}
	W = \frac12 \iint\limits_S \sigma\phi\, dS = \frac12\, \sigma\, \frac{Q}{R}\cdot 4\pi R^2 = \frac{Q^2}{2R}.
\end{equation}

\emphx{Приклад 2.} Нехай тепер той самий заряд $Q$ розподілений рівномірно по \emphx{об'єму} кулі. Обчислення (з урахуванням того, що потенціал
всередині рівномірно зарядженої кулі залежить від $r$) дає інший результат:
\begin{equation}\label{eq:SphereEnergyVolume}
	W = \frac35\, \frac{Q^2}{R}.
\end{equation}

\begin{Attention}
	Той самий сумарний заряд $Q$ на тілі того самого радіуса $R$ дає \emphx{різну} енергію залежно від того, як саме заряд розподілений усередині
	тіла (формули~\eqref{eq:SphereEnergySurface} і~\eqref{eq:SphereEnergyVolume} відрізняються множником $6/5$). Це наштовхує на природне запитання:
	чи можна вважати енергію \emphx{властивістю самих зарядів}, чи, може, правильніше говорити про енергію, \emphx{розподілену в просторі} --- там,
	де існує поле? Відповідь на це запитання ми отримаємо, коли перейдемо до об'ємної густини енергії поля.
\end{Attention}

\begin{Attention}
	\emphx{Класичний радіус електрона.} Цікаво застосувати ці оцінки до електрона. Якщо уявити його як кульку радіуса $r_e$, заряджену рівномірно
	по поверхні, то вся його енергія спокою $m_ec^2$ мала б бути енергією електричного поля:
	\begin{equation*}
		\frac{e^2}{2r_e} = m_ec^2 \quad\Longrightarrow\quad r_e = \frac{e^2}{2m_ec^2}.
	\end{equation*}
	Чисельно це дає $r_e \approx 1{,}4\times 10^{-13}\ \text{см}$ (з точністю до множника порядку одиниці, який залежить від припущення про
	розподіл заряду --- по поверхні чи по об'єму). Цю величину називають \emphx{класичним радіусом електрона}.

	Але тут і криється пастка: якщо електрон справді має такий скінченний розмір, то \emphx{чому} він не розлітається від власного електричного
	відштовхування? І чому досліди з розсіяння на високих енергіях не бачать у нього ніякого розміру аж до $\sim 10^{-16}\ \text{см}$? Відповідь ---
	електрон не є класичною зарядженою кулькою. Класичний радіус --- це не «справжній розмір» електрона, а лише характерний масштаб, на якому
	електрична енергія поля зарядженої частинки стає порівнянною з її енергією спокою. Саме на таких відстанях класична електродинаміка
	остаточно ламається, і на зміну їй приходить квантова теорія поля.
\end{Attention}

%% --------------------------------------------------------
\subsection*{Енергія зарядженого конденсатора}
%% --------------------------------------------------------

Розглянемо процес зарядки конденсатора як послідовне перенесення малих порцій заряду $dq$ з однієї обкладки на іншу (рис.~\ref{pic:CapacitorDq}).
Коли на обкладках уже є заряд $q$ (а різниця потенціалів дорівнює $\Delta\phi = q/C$), перенесення чергової порції $dq$ вимагає роботи
$dW = \Delta\phi\, dq = \dfrac{q\,dq}{C}$. Інтегруючи від нульового заряду до кінцевого значення $q$, знаходимо енергію зарядженого конденсатора:
\begin{equation}\label{eq:CapacitorEnergy}
	\tcbhighmath{
		W = \frac{q^2}{2C} = \frac{C\Delta\phi^2}{2}.
	}
\end{equation}

\begin{SCfigure}[4][h!]\centering
	\localinput{capacitor-dq.tikz}
	\caption{Перенесення заряду $dq$ під час зарядки конденсатора\label{pic:CapacitorDq}}
\end{SCfigure}

\begin{Example}{Числовий приклад} Конденсатор ємністю $C = 1\ \text{мкФ}$, заряджений до напруги $U = 100\ \text{В}$, зберігає заряд
	$q = CU = 10^{-4}\ \text{Кл}$ і енергію $W = CU^2/2 = 5\ \text{мДж}$. Для порівняння: щоб розігнати краплю води масою $1\ \text{мг}$ до
	швидкості $1\ \text{м/с}$, потрібна кінетична енергія $mv^2/2 = 0{,}5\ \text{мкДж}$ --- у \emphx{десять тисяч разів менша}. Саме тому
	конденсатори --- попри скромні, здавалося б, розміри --- здатні зберігати досить значну енергію і широко використовуються як накопичувачі
	(спалахи фотоапаратів, імпульсні джерела живлення, дефібрилятори).
\end{Example}

%% --------------------------------------------------------
\subsection*{Густина енергії електричного поля}\label{sec:FieldEnergyDensity}
%% --------------------------------------------------------

Застосуємо результат~\eqref{eq:CapacitorEnergy} до плоского конденсатора. Оскільки поле в ньому однорідне, $\Delta\phi = Ed$, а ємність дається
формулою~\eqref{eq:FlatCap}, маємо
\begin{equation*}
	W = \frac12 C \Delta\phi^2 = \frac12\,\frac{\epsilon S}{4\pi d}\,(Ed)^2 = \frac{\epsilon E^2}{8\pi}\, V,
\end{equation*}
де $V = Sd$ --- об'єм простору між обкладками, тобто об'єм, що його займає поле. Звідси природно ввести \emphx{об'ємну густину енергії поля}
\begin{equation}\label{eq:CapacitorEnergyDensity}
	\tcbhighmath{
		w = \frac{\epsilon E^2}{8\pi} = \frac{ED}{8\pi}
	}
\end{equation}
(останню рівність отримано підстановкою $D = \epsilon E$; див.~розділ~\ref{FileldInDielectrics}).

\begin{Attention}
	Формула~\eqref{eq:CapacitorEnergyDensity} виражає енергію конденсатора \emphx{не} через заряди на обкладках, а через характеристики поля $\Efield$ і
	$\Dfield$ у тій області простору, де це поле існує. Це дає підставу для важливої фізичної інтерпретації: саме \emphx{електричне поле} є носієм
	взаємодії між зарядами і, водночас, носієм енергії --- енергія \emphx{локалізована} в тих областях простору, де є поле (у плоскому конденсаторі
	--- виключно в проміжку між обкладками), а не \enquote{приписана} безпосередньо зарядам.
\end{Attention}

%% --------------------------------------------------------
\subsection*{Густина енергії електричного поля (строге виведення)}\label{sec:FieldEnergyDensityStrict}
%% --------------------------------------------------------

Формулу~\eqref{eq:CapacitorEnergyDensity} отримано для окремого випадку однорідного поля, локалізованого в проміжку між обкладками плоского конденсатора. Виникає природне запитання: чи має вона більш загальний характер, чи є наслідком саме цієї спрощеної геометрії? Покажемо, що вираз для густини енергії електричного поля залишається справедливим для \emphx{довільного} електростатичного поля --- однорідного чи ні, локалізованого в скінченній області чи створеного як завгодно розподіленими зарядами, --- і виведемо його строго, виходячи з енергії взаємодії зарядів.

Нехай заряд системи змінюється на нескінченно малу величину $\delta\rho(\vect r)$. Оскільки енергія
$W=\tfrac12\iiint\rho\phi\,dV$ квадратична за $\rho$ (адже сам потенціал $\phi$ лінійно виражається
через $\rho$), пряме варіювання дає внесок не лише від зміни $\rho$ при фіксованому $\phi$, а й від
відповідної зміни самого потенціалу. Записуючи $\phi\to\phi+\delta\phi$, де
$\delta\phi(\vect r)=\iiint \delta\rho(\vect r')/|\vect r-\vect r'|\,dV'$ --- потенціал, створений
приростом заряду, --- та користуючись симетрією потенціалу точкового заряду,
\begin{equation*}
	\iiint \rho\,\delta\phi\,dV
	= \iiint\!\!\iiint \frac{\rho(\vect r)\,\delta\rho(\vect r')}{|\vect r-\vect r'|}\,dV\,dV'
	= \iiint \delta\rho\,\phi\,dV,
\end{equation*}
бо відстань між точками не залежить від порядку, $|\vect r-\vect r'|=|\vect r'-\vect r|$. Отже,
\begin{equation}\label{eq:DeltaW}
	\tcbhighmath{
		\delta W = \iiint\limits_V \delta\rho(\vect r)\, \phi(\vect r)\, dV.
	}
\end{equation}

За теоремою Гаусса для вектора $\Dfield$ у диференціальній формі, $\delta\rho = \dfrac1{4\pi}\divg\delta\Dfield$, тому
\begin{equation*}
	\delta W = \frac1{4\pi} \iiint\limits_V \left(\divg \delta\Dfield\right) \phi\, dV.
\end{equation*}
Скористаємось векторною тотожністю $\divg(\phi\, \delta\Dfield) = \phi\, \divg\delta\Dfield + \delta\Dfield\cdot\vect\nabla\phi$, звідки
\begin{equation*}
	\delta W = \frac1{4\pi} \iiint\limits_V \divg\left(\phi\, \delta\Dfield\right) dV - \frac1{4\pi} \iiint\limits_V \delta\Dfield\cdot\vect\nabla\phi\, dV.
\end{equation*}

Поширимо інтегрування на весь (нескінченний) простір, обравши $V$ як велику кулю радіуса $R$ із центром десь усередині системи зарядів, і розглянемо границю $R\to\infty$. Перший доданок за теоремою Остроградського--Гаусса перетворюється на потік крізь сферу $S_R$:
\begin{equation*}
	\oiint\limits_{S_R} \phi\,\delta\Dfield\cdot d\vect S = \oiint\limits_{S_R} \phi\,\delta D_r\, dS.
\end{equation*}
Для обмеженої системи зарядів на великих відстанях $\phi\sim q/R$, а $\delta D_r\sim 1/R^2$, тому підінтегральний вираз спадає як $1/R^3$, тоді як площа сфери зростає як $R^2$. Отже, весь поверхневий інтеграл спадає як $1/R$ і прямує до нуля при $R\to\infty$:
\begin{equation*}
	\oiint\limits_{S_R} \phi\,\delta D_r\, dS \;\sim\; \frac1R \;\xrightarrow[R\to\infty]{}\; 0.
\end{equation*}
Тому перший доданок зникає, і залишається лише
\begin{equation*}
	\delta W = -\frac1{4\pi} \iiint\limits_V \delta\Dfield\cdot\vect\nabla\phi\, dV
	= \iiint\limits_V \frac{\Efield\cdot\delta\Dfield}{4\pi}\, dV,
\end{equation*}
де ми скористались зв'язком $\Efield = -\vect\nabla\phi$.

Для лінійного діелектрика ($\Dfield = \epsilon\Efield$) це співвідношення інтегрується елементарно ($\Efield\cdot\delta\Dfield =
\tfrac12\delta(\Efield\cdot\Dfield)$), і ми приходимо до загального результату:
\begin{equation}\label{eq:GeneralEnergyDensity}
	\tcbhighmath{
		W = \iiint\limits_V \frac{\Efield\cdot\Dfield}{8\pi}\, dV, \qquad w = \frac{\Efield\cdot\Dfield}{8\pi} = \frac{\epsilon E^2}{8\pi}.
	}
\end{equation}

\begin{Attention}
	Формула~\eqref{eq:GeneralEnergyDensity} спирається на лінійний зв'язок $\Dfield=\epsilon\Efield$ з $\epsilon=\text{const}$.
	Саме на кроці $\Efield\cdot\delta\Dfield=\tfrac12\delta(\Efield\cdot\Dfield)$ ми цим скористалися. Проміжний
	результат $\delta w = \Efield\cdot\delta\Dfield/4\pi$, отриманий вище, є при цьому набагато загальнішим: він справедливий для
	\emphx{довільного} середовища --- лінійного чи нелінійного, ізотропного чи ні, --- оскільки спирався лише на скінченність повного
	заряду системи (яка забезпечила зникнення поверхневого інтеграла на нескінченності), а не на конкретний вигляд зв'язку між
	$\Efield$ і $\Dfield$.

	Тому в загальному (нелінійному) випадку густину енергії слід обчислювати як роботу, виконану джерелами поля при поступовому
	наростанні поля від нуля до заданого значення:
	\begin{equation*}
		w = \frac1{4\pi}\int\limits_0^{\Dfield} \Efield(\Dfield')\cdot d\Dfield',
	\end{equation*}
	що для ізотропного середовища ($\Dfield\parallel\Efield$) зводиться до звичайного інтеграла $w=\tfrac1{4\pi}\int_0^{D}E(D')\,dD'$
	--- площі під кривою $E(D)$. Для термодинамічно оборотного (без гістерезису) середовища цей інтеграл не залежить від шляху в
	просторі $\Dfield$, і $w(\Dfield)$ є справжньою функцією стану поля. Легко переконатись, що при $\epsilon=\text{const}$ цей вираз
	відтворює~\eqref{eq:GeneralEnergyDensity}:
	\begin{equation*}
		w = \frac1{4\pi}\int\limits_0^{D} \frac{D'}{\epsilon}\,dD' = \frac{D^2}{8\pi\epsilon} = \frac{\epsilon E^2}{8\pi} = \frac{\Efield\cdot\Dfield}{8\pi}.
	\end{equation*}
\end{Attention}

\begin{Example}{Числовий приклад}
     Оцінімо енергію поля точкового заряду $q$ у вакуумі. Густина енергії $w = E^2/(8\pi) = q^2/(8\pi r^4)$. Інтегруючи
	по сферичних оболонках $dV = 4\pi r^2\, dr$ від деякого малого радіуса $r_0$ до нескінченності:
	\begin{equation*}
		W = \int\limits_{r_0}^{\infty} \frac{q^2}{8\pi r^4}\cdot 4\pi r^2\, dr = \frac{q^2}{2}\int\limits_{r_0}^{\infty} \frac{dr}{r^2} = \frac{q^2}{2r_0}.
	\end{equation*}
	При $r_0 \to 0$ цей інтеграл \emphx{розбігається} --- власна енергія точкового заряду нескінченна. Це --- не фізична катастрофа, а сигнал, що
	\enquote{точковий заряд} --- ідеалізація, яка ламається на дуже малих відстанях (там потрібна квантова теорія поля). Для реальних заряджених
	тіл скінченних розмірів (як куля в прикладах вище) власна енергія залишається скінченною.
\end{Example}

\subsection*{Власна та взаємна енергія зарядів}

Порівняємо два, на перший погляд, різні вирази для енергії системи зарядів: \enquote{дискретну} формулу~\eqref{eq:PairEnergy},
$W = \tfrac12\sum q_iq_j/r_{ij}$, і \enquote{польову} формулу~\eqref{eq:GeneralEnergyDensity}, $W = \int E^2/(8\pi)\, dV$. Перша з них, узагалі кажучи, може
набувати як додатних, так і від'ємних значень (притягання різнойменних зарядів дає від'ємний внесок); друга ж, як інтеграл від невід'ємної величини
$E^2$, --- \emphx{завжди} невід'ємна. Це здається суперечністю, проте насправді обидві формули описують не зовсім те саме.

Розглянемо поле системи двох зарядів, $\Efield = \Efield_1 + \Efield_2$. Тоді
\begin{equation*}
	W = \int\limits_V \frac{E^2}{8\pi}\, dV = \int\limits_V \frac{E_1^2}{8\pi}\, dV + \int\limits_V \frac{E_2^2}{8\pi}\, dV + \int\limits_V
	\frac{\Efield_1\cdot\Efield_2}{8\pi}\, dV.
\end{equation*}
Перші два доданки --- це \emphx{власна енергія} кожного із зарядів окремо; вони завжди додатні,
\begin{equation*}
	\int\limits_V \frac{E_1^2}{8\pi}\, dV > 0, \qquad \int\limits_V \frac{E_2^2}{8\pi}\, dV > 0,
\end{equation*}
і не залежать від взаємного розташування зарядів. Третій доданок --- це \emphx{взаємна (інтерференційна) енергія}
\begin{equation}\label{eq:SelfMutualEnergy}
	W_\text{вз} = \int\limits_V \frac{\Efield_1\cdot\Efield_2}{8\pi}\, dV,
\end{equation}
яка може мати будь-який знак залежно від взаємного розташування й знаків зарядів, і саме вона точно відповідає \enquote{дискретній} формулі
$q_1q_2/r_{12}$ з попереднього пункту.

\begin{Attention}
	Формула~\eqref{eq:PairEnergy} для точкових зарядів описує \emphx{лише взаємну} енергію --- енергію, потрібну для зближення вже готових точкових
	зарядів з нескінченності. Вона свідомо \emphx{не враховує} власну енергію кожного точкового заряду окремо, оскільки поле точкового заряду
	$E \sim 1/r^2$ поблизу нього необмежено зростає, а його власна енергія $\int E^2/(8\pi)\, dV$ --- розбігається (звертається в нескінченність) при
	$r \to 0$. Саме тому \enquote{дискретна} енергія системи точкових зарядів може бути від'ємною --- вона є лише частиною (взаємною частиною) від
	справді повної, завжди додатної, енергії поля, яка додатково містить (нескінченно великі) власні енергії кожного точкового заряду. Для тіл
	скінченних розмірів (як у прикладах з кулею вище) власна енергія, звісно, залишається скінченною.
\end{Attention}

%% --------------------------------------------------------
\section{Пондеромоторні сили}
%% --------------------------------------------------------

Термін \emphx{пондеромоторні сили} походить від латинських слів \textit{pondus} (вага, тягар) та \textit{motor} (той, що рухає). Історично цей термін виник у механіці для позначення сил, які здатні викликати механічний рух тіла. У контексті електродинаміки поняття набуло вужчого значення: \emphx{пондеромоторними} називають механічні сили, що діють на макроскопічні заряджені тіла (провідники, діелектрики, магнітні матеріали) з боку електромагнітного поля.

%% --------------------------------------------------------
\subsection*{Метод віртуальних переміщень}
%% --------------------------------------------------------

Одним з найефективніших способів обчислення сил, що діють на заряджені провідники й діелектрики, є \emphx{метод віртуальних переміщень}. Його ідея
полягає в тому, щоб обчислити роботу сил поля $\delta A = F_x\, dx$ при уявному (\emphx{віртуальному}) нескінченно малому зміщенні тіла на $dx$ у
напрямку, що нас цікавить, після чого шукана сила знаходиться просто як
\begin{equation*}
	F_x = \frac{\delta A}{dx}.
\end{equation*}

\begin{Attention}
	Метод віртуальних переміщень --- універсальний: він не вимагає з'ясування фізичної природи сили (електричної чи механічної) і автоматично
	враховує \emphx{всі} силові взаємодії в системі, оскільки спирається лише на закон збереження енергії.
\end{Attention}

\begin{Attention}
	\emphx{Теорема Томсона.} У методі віртуальних переміщень ми неявно спираємось на глибший факт: рівноважний розподіл заряду на провідниках (за заданих повних зарядів $q_i$ чи заданих потенціалів $\phi_i$) --- \emphx{єдиний}, що мінімізує електростатичну енергію системи серед усіх уявних розподілів заряду, сумісних з накладеними умовами. Інтуїція проста: якби якийсь інший розподіл давав меншу енергію, заряд самочинно перерозподілився б, і рівновага була б неможливою. Це --- своєрідний варіаційний принцип електростатики, тісно пов'язаний з теоремою єдиності розв'язку рівняння Пуассона з попередніх розділів.
\end{Attention}

\subsection*{Сила при сталому заряді провідників}

Нехай на провідниках підтримується сталий заряд, $q = \const$, тобто в системі немає зовнішніх джерел енергії (батарей тощо). За таких умов уся
робота $\delta A$ внутрішніх сил системи при віртуальному переміщенні відбувається виключно за рахунок \emphx{зменшення} власної електричної
енергії системи $W$:
\begin{equation*}
	\delta A = -\left(dW\right)_{q=\const}.
\end{equation*}
Якщо шукана сила $F_x$ здійснює роботу $\delta A = F_x\, dx$ при зміщенні тіла на $dx$, то
\begin{equation}\label{eq:ForceConstQ}
	\tcbhighmath{
		F_x = -\left.\frac{dW}{dx}\right|_{q=\const}.
	}
\end{equation}

\begin{Attention}
	Формула~\eqref{eq:ForceConstQ} --- це суто математичний прийом: для її застосування не обов'язково, щоб заряди провідників у \emphx{реальному}
	процесі залишалися незмінними. Достатньо знайти приріст енергії $dW$, уявляючи, що заряд саме такий, яким він є в даний момент, і \emphx{умовно}
	вважаючи його фіксованим при обчисленні похідної.
\end{Attention}

\subsection*{Сила при сталій різниці потенціалів}

Нехай тепер, навпаки, на провідниках підтримуються сталими потенціали, $\Delta\phi = \const$ (наприклад, провідники приєднані до батареї, яка
підживлює заряди в міру переміщення тіл і виконує при цьому роботу). За законом збереження енергії, робота джерела $\delta A_\text{дж}$ витрачається
одночасно на виконання механічної роботи $F_x\, dx$ і на зміну енергії поля $dW$:
\begin{equation*}
	\delta A_\text{дж} = F_x\, dx + dW.
\end{equation*}
З іншого боку, робота джерела з перенесення заряду $\delta q$ через різницю потенціалів $\Delta\phi$ дорівнює $\delta A_\text{дж} = \delta q\,
\Delta\phi$; оскільки $q = C\Delta\phi$, а $\Delta\phi = \const$, маємо $\delta q = \Delta\phi\, \delta C$, тобто $\delta A_\text{дж} = \Delta\phi^2\,
\delta C$. Водночас, за формулою~\eqref{eq:CapacitorEnergy}, $dW = \tfrac12\Delta\phi^2\, \delta C$ (енергія при сталому $\Delta\phi$ зростає разом
із ємністю). Підставляючи, знаходимо
\begin{equation*}
	F_x\, dx = \Delta\phi^2\,\delta C - \frac12\Delta\phi^2\,\delta C = \frac12\Delta\phi^2\,\delta C = dW,
\end{equation*}
звідки
\begin{equation}\label{eq:ForceConstPhi}
	\tcbhighmath{
		F_x = +\left.\frac{dW}{dx}\right|_{\Delta\phi=\const}.
	}
\end{equation}

\begin{Attention}
	Формули~\eqref{eq:ForceConstQ} і~\eqref{eq:ForceConstPhi} відрізняються \emphx{знаком}, проте описують \emphx{одну й ту саму} фізичну силу: сама
	сила, що діє на нерухоме (у момент обчислення) тіло, не може залежати від того, який уявний спосіб зміни системи ми обираємо для її обчислення
	--- метод сталого заряду чи метод сталого потенціалу. Різниця знаків пояснюється тим, що при сталому $\Delta\phi$ енергію в систему додатково
	вкладає джерело, і саме цей додатковий внесок \enquote{перевертає} знак похідної.

	\emphx{Мнемонічне правило:} якщо система \enquote{віддана сама собі} ($q = \const$) --- вона \enquote{прагне} зменшити свою енергію; якщо
	система підживлюється джерелом ($\Delta\phi = \const$) --- джерело \enquote{допомагає} силі, і знак стає додатним. Але фізична сила ---
	однакова в обох випадках.
\end{Attention}


\begin{Example}{Притягування обкладок зарядженого конденсатора}

Знайдемо силу взаємного притягання обкладок плоского конденсатора, уявивши, що одна з обкладок зазнає віртуального зміщення $dx$ вздовж напрямку,
що з'єднує обкладки (рис.~\ref{pic:VirtualDisplacement}).

\begin{center}%[2][h!]\centering
	\localinput{virtual.tikz}
	\captionof{figure}{Віртуальне зміщення $dx$ обкладки конденсатора\label{pic:VirtualDisplacement}}
\end{center}

\emphx{Спосіб 1: $q = \const$.} Зміна енергії поля при зміні відстані між обкладками:
\begin{equation*}
	dW = d\left(\frac{q^2}{2C}\right) = -\frac{q^2}{2C^2}\, dC, \qquad C = \frac{\epsilon S}{4\pi x}, \qquad dC = -\frac{C}{x}\, dx,
\end{equation*}
звідки $dW = \dfrac{q^2}{2Cx}\, dx$. Виражаючи через напруженість поля ($q = CEd = C E x$, де $x$ тут відіграє роль відстані між обкладками),
\begin{equation*}
	dW = \frac{E^2 x^2\, \epsilon S}{2 \cdot 4\pi x \cdot x}\, dx = \frac{\epsilon E^2 S}{8\pi}\, dx,
\end{equation*}
і за формулою~\eqref{eq:ForceConstQ}
\begin{equation}\label{eq:ForceCapacitorPlates}
		F_x = -\frac{dW}{dx} = -\frac{\epsilon E^2}{8\pi}\, S.
\end{equation}

\emphx{Спосіб 2: $\Delta\phi = \const$.} Аналогічно, $dW = d\left(\dfrac{C\Delta\phi^2}{2}\right) = \dfrac{\Delta\phi^2}{2}\, dC =
-\dfrac{\Delta\phi^2 C}{2x}\, dx$, що після підстановки $\Delta\phi = Ex$ дає $dW = -\dfrac{\epsilon E^2 S}{8\pi}\, dx$, і за
формулою~\eqref{eq:ForceConstPhi} знову отримуємо той самий результат~\eqref{eq:ForceCapacitorPlates}.


Знак \enquote{мінус} у формулі~\eqref{eq:ForceCapacitorPlates} свідчить про те, що сила спрямована на \emphx{зменшення} відстані між обкладками, тобто
обкладки \emphx{притягуються} одна до одної --- незалежно від знаків зарядів на них, оскільки сила пропорційна $E^2$.
\end{Example}

\begin{Attention}
	Результат~\eqref{eq:ForceCapacitorPlates} має набагато загальніший характер, ніж просто «сила між обкладками конденсатора». Тиск поля на одиницю площі будь-якої зарядженої поверхні провідника дорівнює
	\[
		p = \frac{E^2}{8\pi} = 2\pi\sigma^2,
	\]
	де $E$ --- напруженість поля безпосередньо біля поверхні (зовні провідника), а $\sigma$ --- поверхнева густина заряду в цій точці (тут використано $E=4\pi\sigma$ біля поверхні провідника). Це --- \emphx{електростатичний тиск}, що намагається «розірвати» будь-яке заряджене тіло; він діє в кожній точці поверхні незалежно від її форми і завжди спрямований \emphx{назовні} провідника. У випадку плоского конденсатора цей тиск, помножений на площу обкладки $S$, і дає силу~\eqref{eq:ForceCapacitorPlates}.
\end{Attention}

\begin{Example}{Сила втягування діелектрика в конденсатор}
Розглянемо плоский конденсатор, у який частково (на глибину $x$) всунуто пластину діелектрика з проникністю $\epsilon$, і знайдемо силу, що діє на
цю пластину. Зручно замінити реальну систему еквівалентною --- двома паралельно з'єднаними конденсаторами: один, площею $S_1$, --- вакуумний, а
інший, площею $S_2$, --- повністю заповнений діелектриком.

\begin{center}%[h!]\centering
     \includegraphics[height=7cm]{pull_dielectric}
     \captionof{figure}{Модель конденсатора з частково всунутою пластиною діелектрика: паралельне з'єднання вакуумного та заповненого діелектриком конденсаторів\label{pic:DielectricInCapacitor}}
\end{center}

За формулою для паралельного з'єднання~\eqref{eq:ParallelCap},
\begin{equation*}
	C = C_1 + C_2 = \frac{S_1}{4\pi d} + \frac{\epsilon S_2}{4\pi d}.
\end{equation*}
Нехай повна довжина конденсатора (у напрямку висування) дорівнює $L$, ширина обкладок --- $h$, тоді повна площа $S = Lh$ стала, а
$S_1 = \left(1 - \dfrac{x}{L}\right) S$, $S_2 = \dfrac{x}{L}S$ (тут $x$ --- глибина занурення пластини діелектрика). Підставляючи, знаходимо ємність
як функцію $x$:
\begin{equation*}
	C(x) = \frac{S}{4\pi d}\left(1 + (\epsilon - 1)\frac{x}{L}\right).
\end{equation*}
Обчислимо силу при сталому заряді $q = \const$ за формулою~\eqref{eq:ForceConstQ}:
\begin{equation*}
	F_x = -\frac{dW}{dx} = -\frac{dW}{dC}\frac{dC}{dx}, \qquad \frac{dW}{dC} = -\frac{q^2}{2C^2}, \qquad \frac{dC}{dx} = \frac{S}{4\pi d}\,
	\frac{\epsilon-1}{L},
\end{equation*}
звідки
\begin{equation*}
	F_x = \frac{q^2}{2C^2}\cdot\frac{S}{4\pi d}\cdot\frac{\epsilon - 1}{L} \stackrel{q = CEd}{=} \frac{(\epsilon-1)E^2}{8\pi}\, hd,
\end{equation*}
де ми скористались тим, що в кожен момент $q = CEd$, а $S/L = h$. Остаточно,
\begin{equation}\label{eq:ForceDielectric}
		F_x = \frac{(\epsilon-1)E^2}{8\pi}\, S_\text{торця},
\end{equation}
де $S_\text{торця} = hd$ --- площа торцевого перерізу пластини діелектрика (перерізу, яким вона \enquote{входить} у конденсатор). Оскільки $\epsilon > 1$,
сила $F_x$ додатна --- вона спрямована в бік \emphx{втягування} діелектрика всередину конденсатора.
\end{Example}

\begin{Attention}
	Результат~\eqref{eq:ForceDielectric} має простий фізичний зміст: діелектрик завжди \emphx{втягується} в область сильнішого поля, оскільки
	заповнення проміжку поляризовним середовищем енергетично вигідне --- воно збільшує ємність (при сталому заряді енергія при цьому \emphx{падає}),
	а система, як завжди в природі, \enquote{прагне} до стану з меншою енергією. Це той самий фізичний механізм, завдяки якому рідкий діелектрик
	підіймається між частково зануреними в нього обкладками зарядженого конденсатора.
\end{Attention}

\begin{Attention}
	\emphx{Зв'язок із діелектриками.} Зверніть увагу: у формулі~\eqref{eq:ForceDielectric} сила пропорційна $(\epsilon - 1)$ --- тобто вона
	\emphx{дорівнює нулю} для вакууму ($\epsilon = 1$) і зростає зі збільшенням поляризовності середовища. Це природно: сила втягування --- це
	прямий прояв \emphx{поляризації} діелектрика в неоднорідному полі, яку ми розбирали в розділі~\ref{FileldInDielectrics}. Чим сильніше
	середовище поляризується, тим сильніше воно \enquote{притягується} до області сильного поля.
\end{Attention}

%% --------------------------------------------------------
\section*{Підсумки}
%% --------------------------------------------------------

Усі результати цього розділу спираються на одну спільну ідею: електричне поле --- не допоміжна розрахункова абстракція, а фізичний об'єкт, який
\emphx{сам є носієм енергії}. Ємність описує, скільки заряду (а отже, і поля) система здатна накопичити при заданій різниці потенціалів; енергія
поля визначається інтегралом густини $w = \Efield\cdot\Dfield/8\pi$ за областю, де це поле існує; а пондеромоторні сили, що діють на провідники й
діелектрики, --- це прояв фундаментального прагнення будь-якої фізичної системи до мінімуму енергії.

Практично з цього розділу варто винести три речі. По-перше, ємність будь-якої системи провідників обчислюється за єдиним алгоритмом: задати заряд,
знайти поле, проінтегрувати різницю потенціалів, поділити. По-друге, густина енергії поля $w = \Efield\cdot\Dfield/8\pi$ --- формула загальна, а
не лише для плоского конденсатора: вона справедлива для довільного електростатичного поля. По-третє, сили, що діють на провідники й діелектрики,
найпростіше знаходити методом віртуальних переміщень --- через похідну енергії за координатою, з урахуванням того, яка величина (заряд чи
потенціал) підтримується сталою.

%% --------------------------------------------------------
\section*{Задачі}
%% --------------------------------------------------------

\begin{problem}\label{prob:SphereEnergy}
	\emphx{Енергія поля зарядженої кулі.} Обчисліть енергію електричного поля рівномірно зарядженої по об'єму кулі радіуса $R$ із повним
	зарядом $Q$ двома способами: (а) через формулу~\eqref{eq:SphereEnergyVolume}; (б) через інтеграл $\int E^2/(8\pi)\, dV$ за всім простором.
	Переконайтесь, що результати збігаються. \emphx{Вказівка:} поле всередині кулі $E = Qr/R^3$, зовні --- $E = Q/r^2$.
\end{problem}

\begin{problem}\label{prob:DielectricForce}
	\emphx{Сила, що діє на діелектрик.} У плоский конденсатор із площею обкладок $S = 50\ \text{см}^2$ і відстанню $d = 2\ \text{мм}$ частково
	всунуто пластину діелектрика з $\epsilon = 4$. Поле в конденсаторі $E = 5\times 10^4\ \text{В/м}$. Обчисліть силу, що втягує пластину
	всередину конденсатора. \emphx{Відповідь:} $F \approx 1{,}1\times 10^{-3}\ \text{Н}$.
\end{problem}

\begin{problem}\label{prob:CapacitorBank}
	\emphx{Ємність батареї.} Три конденсатори ємностями $C_1 = 1\ \text{мкФ}$, $C_2 = 2\ \text{мкФ}$, $C_3 = 3\ \text{мкФ}$ з'єднано
	спочатку послідовно, потім паралельно. Обчисліть сумарну ємність у кожному випадку. \emphx{Відповідь:} послідовно $\approx 0{,}55\ \text{мкФ}$;
	паралельно $6\ \text{мкФ}$.
\end{problem}

\begin{problem}\label{prob:CapacitanceMatrix}
	\emphx{Матриця ємностей.} Для двох куль із попереднього прикладу покладіть $R_1=1$~см, $R_2=2$~см, $b=50$~см і обчисліть
	$C_{11}, C_{22}, C_{12}$ чисельно. \emphx{Відповідь:} $C_{11} \approx 1$~см, $C_{22} \approx 2$~см, $C_{12} \approx -0{,}04$~см (знак від'ємний, оскільки $C_{12} \approx -R_1R_2/b$).
\end{problem}

\begin{problem}\label{prob:ForceConstPhi}
	\emphx{Сила при сталій напрузі.} Плоский конденсатор із площею обкладок $S=20\ \text{см}^2$ приєднано до батареї, що підтримує напругу $U=300$~В; відстань між обкладками $d=0{,}5$~мм. Обчисліть силу притягання обкладок безпосередньо за формулою~\eqref{eq:ForceConstPhi} (через залежність $C(x)$), не користуючись готовим результатом~\eqref{eq:ForceCapacitorPlates}, і переконайтесь, що відповідь збігається.
\end{problem}

\begin{Questions}
\begin{quizlist}
	\item Що називають ємністю усамітненого провідника? Чому відношення $C = q/\phi$ не залежить від заряду і від чого воно залежить насправді?
	\item Чому означення $C = q/\phi$ втрачає сенс, якщо поблизу є інші провідники? Як тоді визначається взаємна ємність пари провідників?
	\item Що таке конденсатор і чому його взаємна ємність майже не залежить від оточуючих тіл?
	\item Чому дорівнює ємність двох далеких одна від одної куль (рис.~\ref{pic:TwoSpheres}) у головному наближенні? Що описує поправка, пропорційна $1/b$?
	\item Виведіть формули для ємності батареї при послідовному та паралельному з'єднанні конденсаторів (рис.~\ref{pic:CapacitorsConnection}). Що залишається однаковим у кожному з випадків: заряд чи напруга?
	\item Чому послідовне з'єднання завжди зменшує сумарну ємність, а паралельне збільшує? У чому відмінність від з'єднання резисторів?
	\item Сформулюйте триетапний алгоритм обчислення ємності. Застосуйте його для плоского конденсатора та запишіть результат.
	\item Виведіть ємність сферичного конденсатора. Що станеться з формулою при $R_2\to\infty$ і чому цей результат очікуваний?
	\item Яка загальна структура формул для ємності плоского, сферичного та циліндричного конденсаторів? Що в них відіграє роль «геометричного розміру»?
	\item Чому дорівнює взаємна енергія системи точкових зарядів? Навіщо в формулі $U = \tfrac12\sum\sum q_iq_j/r_{ij}$ множник $\tfrac12$ і чому вона не враховує власну енергію зарядів?
	\item Запишіть енергію системи зарядів через потенціали. Чому енергія кулі, зарядженої по поверхні, відрізняється від енергії кулі того ж заряду й радіуса, зарядженої по об'єму? Яке запитання це ставить про локалізацію енергії?
	\item Виведіть енергію зарядженого конденсатора $W = q^2/2C = C\Delta\phi^2/2$ через роботу перенесення заряду $dq$.
	\item Як з енергії плоского конденсатора отримують густину енергії поля $w = \epsilon E^2/8\pi = ED/8\pi$? Чому цю формулу можна вважати загальною, а не лише для однорідного поля?
	\item Яку роль відіграє те, що потік вектора $\phi\,\delta\Dfield$ крізь нескінченно віддалену поверхню дорівнює нулю, у виведенні загального виразу для густини енергії?
	\item Чому енергія «дискретної» формули для точкових зарядів може бути від'ємною, а енергія, обчислена як інтеграл від $E^2/8\pi$, завжди додатна? Що таке власна та взаємна (інтерференційна) енергія?
	\item Чому власна енергія точкового заряду нескінченна? Що це означає для моделі «точковий заряд»?
	\item Що таке пондеромоторні сили та звідки походить термін? У чому полягає метод віртуальних переміщень?
	\item Виведіть формули для сили при $q=\const$ та при $\Delta\phi=\const$. Чому вони відрізняються знаком, хоча описують одну й ту саму фізичну силу?
	\item Знайдіть силу притягання обкладок плоского конденсатора. Чому вона напрямлена на зменшення відстані між обкладками незалежно від знаків зарядів?
	\item Чому діелектрик втягується в область сильного поля конденсатора? Як у формулі для сили з'являється множник $(\epsilon-1)$ і що з цього випливає для вакууму?
\end{quizlist}
\end{Questions}